\renewcommand{\bibname}{Используемая литература}

\begin{thebibliography}{99}
\addcontentsline{toc}{chapter}{Используемая литература}

\bibitem{kernighan-ritchie-c}
Керниган Б., Ритчи Д. \textit{Язык программирования C}. Классический источник по языку C и его модели программирования.

\bibitem{bryant-ohallaron-csapp}
Bryant R. E., O'Hallaron D. R. \textit{Computer Systems: A Programmer's Perspective}. Базовый источник о связи программ, памяти, процессора, компиляции и операционной системы.

\bibitem{patterson-hennessy}
Patterson D. A., Hennessy J. L. \textit{Computer Organization and Design}. Введение в архитектуру компьютеров и машинное представление данных.

\bibitem{intel-manual}
Intel. \textit{Intel 64 and IA-32 Architectures Software Developer's Manual}. Справочная документация по архитектуре x86-64.

\bibitem{nasm-doc}
The NASM Development Team. \textit{NASM Documentation}. Справочная документация по ассемблеру NASM.

\bibitem{gcc-manual}
Free Software Foundation. \textit{Using the GNU Compiler Collection}. Документация по GCC и этапам компиляции.

\bibitem{gnu-make-manual}
Free Software Foundation. \textit{GNU Make Manual}. Документация по системе сборки \texttt{make}.

\bibitem{pro-git}
Chacon S., Straub B. \textit{Pro Git}. Практическое введение в Git и командную разработку.

\bibitem{linux-man-pages}
The Linux man-pages project. \textit{Linux Programmer's Manual}. Справочные страницы Linux для системных вызовов, библиотечных функций и команд.

\bibitem{goldberg-floating-point}
Goldberg D. \textit{What Every Computer Scientist Should Know About Floating-Point Arithmetic}. Классическая статья о числах с плавающей точкой и ошибках округления.

\bibitem{ieee-754}
IEEE. \textit{IEEE Standard for Floating-Point Arithmetic}. Стандарт представления и арифметики чисел с плавающей точкой.

\end{thebibliography}
